\chapter{3D Acquisition}\label{ch:acq}

Every pipeline in this course starts with measurements of a real object. This chapter reviews how
3D scanners work, what their measurements look like, and why the raw output (several noisy,
partial point clouds in different coordinate systems) needs the registration and reconstruction
steps of the next chapters. It follows the first part of Tema 5; the taxonomy follows
\citet{Hartley2004} for stereo and the survey of structured light by \citet{Geng2011}.

\section{From images to 3D data}

A digital camera samples the light that reaches a sensor, producing a 2D array of intensities. A
\emph{range} sensor produces, for each pixel, a distance along the corresponding viewing ray, a
\emph{range image}. Knowing the intrinsic calibration of the sensor, each pixel $(u,v)$ with depth
$Z$ back-projects to a 3D point
\begin{equation}
\bp=Z\,\mK^{-1}\begin{pmatrix}u\\v\\1\end{pmatrix},\qquad
\mK=\begin{pmatrix}f_x&0&c_x\\0&f_y&c_y\\0&0&1\end{pmatrix},
\end{equation}
so a range image is an organised point cloud in the coordinate system of the sensor. Points with
an estimated normal are sometimes called \emph{surfels}. Point clouds are what every capture device
produces, which is the first reason to study them; the second is that they are local and simple.
Their limitations are the absence of connectivity, of topology and of a smooth rendering, together
with noise.

\section{A taxonomy of scanners}

\begin{itemize}
  \item \textbf{Contact} devices (coordinate measuring machines, articulated arms) touch the object
  with a probe: very accurate, slow, and unsuitable for fragile objects.
  \item \textbf{Transmissive} devices (CT, MRI) measure the interior and produce volumes.
  \item \textbf{Optical, active} devices emit light and observe its return: time-of-flight and
  phase-based range scanners, laser triangulation, structured light, active stereo.
  \item \textbf{Optical, passive} methods use only ambient light: stereo, multi-view stereo and
  photogrammetry, shape from shading.
\end{itemize}
The course focuses on single-view optical scanners, which measure one side of the object at a time,
grouped in three families: range scanners, triangulation scanners and stereo.

\section{Time-of-flight scanners}

A pulsed time-of-flight scanner emits a short laser pulse, detects its reflection and measures the
round-trip time $T$. Since light travels to the surface and back,
\begin{equation}
D=\frac{c\,T}{2},\qquad \Delta D=\frac{c\,\Delta T}{2},
\end{equation}
where $c\approx3\times10^8$\,m/s. The depth resolution is limited by the timing resolution: a
1\,GHz clock ($\Delta T=1$\,ns) gives $\Delta D=15$\,cm, and millimetre resolution requires
picosecond timing. Pulsed scanners therefore suit long ranges (buildings, terrain, autonomous
vehicles) where centimetre accuracy suffices. A LiDAR sweeps the beam with rotating mirrors or a
spinning head to cover the scene.

\section{Phase-based scanners}

A continuous-wave scanner emits light whose intensity is modulated at frequency $f$,
$s(t)=a\cos(2\pi ft)$, and receives an attenuated, delayed copy plus ambient light,
$r(t)=A\cos(2\pi f(t-\tau))+B$, with $\tau=2D/c$. The delay appears as a phase shift
$\varphi=2\pi f\tau$, measured with the cross-correlation
\begin{equation}
C(x)=\lim_{T\to\infty}\frac1T\int_{-T/2}^{T/2}s(t)\,r(t-x)\dd t=\frac{aA}{2}\cos\bigl(2\pi f(x+\tau)\bigr)=\frac{aA}{2}\cos(2\pi fx+\varphi) .
\end{equation}
Sampling $C$ at four offsets $x_k=k/(4f)$, $k=0,\dots,3$ (the ``four-bucket'' method) gives
$C_0=K\cos\varphi$, $C_1=-K\sin\varphi$, $C_2=-K\cos\varphi$ and $C_3=K\sin\varphi$ with
$K=aA/2$, so
\begin{equation}
\varphi=\operatorname{atan2}(C_3-C_1,\;C_0-C_2),\qquad A\propto\sqrt{(C_3-C_1)^2+(C_0-C_2)^2},\qquad
D=\frac{c\,\varphi}{4\pi f} .
\end{equation}
The ambient term $B$ cancels in the differences, and the amplitude measures the confidence of the
pixel. Because the phase is only known modulo $2\pi$, distances are ambiguous beyond
$D_{\max}=c/(2f)$ (7.5\,m for $f=20$\,MHz); using two modulation frequencies resolves the
ambiguity. Phase-based scanners are more accurate than pulsed ones at short and medium range, and
time-of-flight cameras (such as the second Kinect) apply the same principle to every pixel of a
sensor. In general, range scanners provide limited precision and their output requires substantial
post-processing.

\fig[0.85\textwidth]{ch06_phase.pdf}{Phase-based ranging. Left: emitted and received signals; the
received one is attenuated, offset by ambient light and delayed by $\tau=2D/c$. Right: the
cross-correlation is a sinusoid whose phase $\varphi=2\pi f\tau$ is recovered from four samples
(the four buckets).}{fig:phase}

\section{Triangulation scanners}

A triangulation scanner adds a camera to a light emitter at a known distance $b$, the
\emph{baseline}. The emitter projects a laser dot or line in a known direction, and the camera
observes where it lands on the object. The emitter, the camera and the illuminated point form a
triangle with one known side ($b$) and two known angles: $\alpha$ at the emitter, from the
projection direction, and $\beta$ at the camera, from the pixel where the dot appears. By the law
of sines,
\begin{equation}
d=\frac{b\sin\alpha}{\sin(\alpha+\beta)},\qquad Z=d\sin\beta=\frac{b\sin\alpha\sin\beta}{\sin(\alpha+\beta)},
\end{equation}
where $d$ is the distance from the camera to the point and $Z$ its distance to the baseline
(\cref{fig:triangulation}). Equivalently, the shift of the observed dot with respect to its
position on a reference plane encodes the depth. Well-calibrated laser triangulation reaches
extreme precision (tens of micrometres), which is why it was used for the Digital Michelangelo
project \citep{Levoy2000}, but sweeping a single line over the object is slow and the equipment is
expensive. Depth precision degrades with distance, and the occlusions seen by either the emitter or
the camera create holes.

\begin{figure}[tbp]
\centering
\begin{tikzpicture}[scale=1.0, every node/.style={font=\small}]
  \coordinate (L) at (0,0); \coordinate (Cm) at (4,0); \coordinate (P) at (1.6,3.0);
  \draw[thick] (L) -- node[below] {baseline $b$} (Cm);
  \draw[BrickRed, thick, -{Stealth}] (L) -- (P);
  \draw[MidnightBlue, thick] (P) -- (Cm);
  \draw[dashed, gray] (P) -- (1.6,0) node[below, black] {};
  \node[left] at ($(L)!0.5!(P)$) {laser};
  \node[right] at ($(Cm)!0.5!(P)$) {$d$};
  \node at (1.85,1.4) {$Z$};
  \fill (P) circle (2pt) node[above] {$\bp$};
  \draw (0.55,0) arc (0:62:0.55); \node at (0.75,0.32) {$\alpha$};
  \draw (3.45,0) arc (180:118:0.55); \node at (3.25,0.32) {$\beta$};
  \draw[fill=black!10] (L) ++(-0.25,-0.2) rectangle ++(0.5,0.4);
  \draw[fill=MidnightBlue!15] (Cm) ++(-0.3,-0.22) rectangle ++(0.6,0.44);
  \node at ($(Cm)+(0.15,-0.55)$) {camera};
  % stereo sketch
  \begin{scope}[xshift=7.9cm]
    \coordinate (O1) at (0,0); \coordinate (O2) at (3,0); \coordinate (Q) at (1.2,3.0);
    \draw[thick] (O1) -- node[below] {$B$} (O2);
    \draw (-0.6,0.6) -- (0.6,0.6); \draw (2.4,0.6) -- (3.6,0.6);
    \node[left] at (-0.6,0.6) {image planes};
    \draw[MidnightBlue] (O1) -- (Q); \draw[MidnightBlue] (O2) -- (Q);
    \fill (Q) circle (2pt) node[above] {$\bp$};
    \fill[BrickRed] (0.24,0.6) circle (1.5pt); \fill[BrickRed] (2.64,0.6) circle (1.5pt);
    \draw[<->] (0,0.85) -- node[above, font=\scriptsize] {$x_l$} (0.24,0.85);
    \draw[<->] (2.64,0.85) -- node[above, font=\scriptsize] {$x_r$} (3,0.85);
    \draw[<->, gray] (3.9,0) -- node[right] {$f$} (3.9,0.6);
    \draw[<->, gray] (4.5,0) -- node[right] {$Z$} (4.5,3.0);
    \node at (1.5,-0.75) {$d=x_l-x_r=fB/Z$};
  \end{scope}
\end{tikzpicture}
\caption{Left: laser triangulation; the depth follows from the baseline and the two angles. Right:
rectified stereo; a point at depth $Z$ appears in the two images with disparity $d=fB/Z$.}
\label{fig:triangulation}
\end{figure}

\section{Structured light}

Structured-light scanners replace the laser by a projector, which illuminates the whole scene with
patterns. The geometry is that of triangulation, with the projector acting as an inverse camera;
the challenge is to identify, for each camera pixel, which projector column (stripe) illuminated it.
Projecting a single pattern of $N$ stripes is ambiguous, but a sequence of binary patterns gives
each stripe a code: with $n$ patterns, $2^n$ stripes can be distinguished, so
$\lceil\log_2N\rceil$ projections suffice for $N$ stripes. The natural binary code has the defect
that many bits change between neighbouring stripes, so a decoding error at a stripe boundary can
produce a large error. The reflected binary \emph{Gray code} changes exactly one bit between
consecutive stripes, so boundary errors cost at most one stripe \citep{Inokuchi1984}; the Gray
code of $k$ is $k\oplus(k\gg1)$ (\cref{fig:graycode}). Sub-stripe accuracy is obtained by
additionally projecting shifted sinusoids and measuring their phase per pixel, exactly as in the
phase-based scanner.

\fig[0.85\textwidth]{ch06_graycode.pdf}{Binary and Gray code patterns for $2^4=16$ stripes, one row
per projected pattern. In the Gray code consecutive columns differ in exactly one bit, so an error
at a stripe boundary displaces the decoded stripe by at most one.}{fig:graycode}

Structured light is fast and cheap (a projector and a camera), and acquires a dense range image
in a fraction of a second. It needs controlled conditions, since strong ambient light washes out
the patterns, and a careful calibration of the camera--projector pair. The first Microsoft
Kinect (2009--2010) made the idea a consumer product: an infrared projector casts a fixed
pseudo-random dot pattern that is matched in a single shot against a stored reference, producing
$640\times480$ depth images (about $300\,000$ points) at 30 frames per second, with precision from
about 1\,mm at 0.5\,m to 4\,cm at 2\,m.

\section{Stereo}

Passive stereo uses two calibrated cameras. After \emph{rectification}, which makes the image
rows correspond (epipolar lines horizontal), a point at depth $Z$ appears at horizontal
coordinates $x_l$ and $x_r$ in the two images, and similar triangles give the \emph{disparity}
\begin{equation}
d=x_l-x_r=\frac{fB}{Z},
\end{equation}
where $f$ is the focal length in pixels and $B$ the baseline (\cref{fig:triangulation}). Inverting,
$Z=fB/d$, and differentiating,
\begin{equation}
\Delta Z\approx\frac{Z^2}{fB}\,\Delta d :
\end{equation}
the depth error grows quadratically with distance and decreases with the baseline and the focal
length (\cref{fig:stereo}). The hard part of stereo is not the geometry but the \emph{correspondence
problem}: finding, for each pixel of the left image, the pixel of the right image that sees the
same point. It fails on textureless or repetitive regions, which is why \emph{active stereo}
projects a random texture to create matchable features. With many images from unknown positions,
structure from motion recovers the cameras and multi-view stereo the dense geometry
(photogrammetry).

\fig[0.85\textwidth]{ch06_stereo.pdf}{Stereo geometry for a Kinect-like rig ($f=580$\,px,
$B=7.5$\,cm). Left: disparity against depth, $d=fB/Z$. Right: depth error for a disparity error of
$1/8$ pixel; it grows as $Z^2$, from about 3\,mm at 1\,m to 46\,mm at 4\,m.}{fig:stereo}

\section{Modern devices}

Today's sensors combine these principles: solid-state and rotating LiDARs (pulsed time of flight)
for robotics and mapping, time-of-flight cameras and LiDAR in mobile phones, structured light for
face authentication, active stereo cameras for robotics, and photogrammetry from ordinary photos.
Learning-based monocular depth estimation complements them when only one image is available.

\section{What a scan looks like}

All of these devices deliver data with the same defects, which the rest of the course must handle:
\begin{itemize}
  \item \textbf{Partiality}: a single view sees only one side of the object; a complete model
  requires several scans.
  \item \textbf{Different coordinate systems}: each scan is expressed in the frame of the sensor at
  the moment of acquisition (\cref{fig:scans}).
  \item \textbf{Noise}: depth errors grow with distance and with the angle between the viewing ray
  and the surface.
  \item \textbf{Outliers}: spurious points at depth discontinuities (mixed pixels), reflections and
  dust.
  \item \textbf{Holes}: occlusions, and dark, transparent or specular materials that return no
  signal.
\end{itemize}
The processing pipeline that turns scans into a model is therefore: acquire the scans, align them
(registration, \cref{ch:icp,ch:global}), remove noise and outliers, estimate normals, and reconstruct
a surface (\cref{ch:recon}).

\fig[0.85\textwidth]{ch06_scans.png}{Two simulated range scans of the bunny from different
viewpoints; each covers only part of the surface. Overlaid in the world frame in which they were
simulated, they complement each other. A real scanner delivers each scan in its own sensor frame,
and registration must recover the rigid transformations that bring them into a common
frame.}{fig:scans}

\begin{keypoints}
\begin{itemize}[leftmargin=*]
  \item Pulsed time of flight: $D=cT/2$; long range, resolution limited by timing (1\,ns gives
  15\,cm).
  \item Phase-based ranging: $D=c\varphi/(4\pi f)$, recovered with four correlation samples,
  ambiguous beyond $c/(2f)$.
  \item Triangulation and structured light: depth from a known baseline and two angles; Gray codes
  identify $N$ stripes with $\lceil\log_2N\rceil$ patterns.
  \item Stereo: $d=fB/Z$, error growing as $Z^2$; the difficulty is finding correspondences.
  \item Scans are partial, noisy, have outliers and holes, and live in different coordinate systems.
\end{itemize}
\end{keypoints}

\section*{Exercises}
\begin{exercise}
A phase-based scanner uses $f=10$\,MHz. What is its unambiguous range? If the phase is measured
with an error of $0.01$\,rad, what is the distance error?
\end{exercise}
\begin{exercise}
Write the Gray code of the integers $0,\dots,7$ and check that consecutive codes differ in one bit.
How do you decode a Gray code back to the integer?
\end{exercise}
\begin{exercise}
A stereo rig has $f=800$\,px and $B=0.1$\,m, and matches with $0.25$\,px accuracy. Compute the
depth error at $1$\,m and at $5$\,m.
\end{exercise}
